% GENERATED FILE. Edit canonical lessons, then run npm run generate:books.
% canonical-source-hash: fnv1a64:7a705d69d8686384
% canonical-lessons: MR-C544-arakhda, MR-C544-masuda, MR-C544-niyojan-karne, MR-C544-lambnivar-takne, MR-C544-tadis-nene

\chapter{Outlines and Drafts}
\label{ch:mr-outlines-and-drafts}
\hlchaptermodality{544}

\begin{chapteropening}
\textbf{By the end of this chapter:} \emph{I can talk about outlines, drafts and seeing plans through.}

Complete the last lesson of chapter 544: I can talk about outlines, drafts and seeing plans through.
\end{chapteropening}

\section[\texorpdfstring{\mr{आराखडा} (ārākhḍā)}{आराखडा (ārākhḍā)}]{\mr{आराखडा} (ārākhḍā) — an outline; a layout}

\label{lesson:MR-C544-arakhda}



\begin{quote}
Before the new one: say the Marathi for valid, then the Marathi for an identity card.
\end{quote}

\subsection*{You'll want to know: \mr{आराखडा}}
\textbf{\mr{आराखडा}} — \emph{ārākhḍā} — \textquotedblleft{}an outline; a layout\textquotedblright{}.

\begin{quote}
\mr{अभियंत्याने} \mr{नवीन} \mr{इमारतीचा} \mr{आराखडा} \mr{बनवला}.
\end{quote}

\begin{quote}
\emph{abhiyantyāne navīn imārtīcā ārākhḍā banavlā.}
\end{quote}

\begin{quote}
The engineer made the layout of the new building.
\end{quote}

Masculine.

\begin{cousinweb}
\textbf{\mr{आराखडा}} (\emph{ārākhḍā}) is a Marathi word for an outline plan; its origin is not settled.
\end{cousinweb}

\begin{grammarlens}[title={What you've built}]
One more word for saying what happens if.
\end{grammarlens}

\subsection*{Your turn}
\begin{itemize}
\raggedright
  \item \emph{Say it:} \emph{ārākhḍā}
  \item \emph{Say it:} \emph{ārākhḍā}, once more
  \item \emph{Say it:} \emph{abhiyantyāne navīn imārtīcā ārākhḍā banavlā.} — The engineer made the layout of the new building.
  \item \emph{Recall:} say the Marathi for valid, then the Marathi for an identity card, then say \emph{ārākhḍā} again
\end{itemize}

\begin{tcolorbox}[breakable,colback=teal!4,colframe=teal!35!black,title={Before you move on}]
What does \mr{आराखडा} mean? (\textquotedblleft{}An outline; a layout\textquotedblright{}.) Say it once more.
\end{tcolorbox}
\section[\texorpdfstring{\mr{मसुदा} (masudā)}{मसुदा (masudā)}]{\mr{मसुदा} (masudā) — a draft}

\label{lesson:MR-C544-masuda}



\begin{quote}
Before the new one: say the Marathi for an identity card, then the Marathi for an outline.
\end{quote}

\subsection*{You'll want to know: \mr{मसुदा}}
\textbf{\mr{मसुदा}} — \emph{masudā} — \textquotedblleft{}a draft\textquotedblright{}.

\begin{quote}
\mr{पत्राचा} \mr{पहिला} \mr{मसुदा} \mr{तयार} \mr{आहे}.
\end{quote}

\begin{quote}
\emph{patrācā pahilā masudā tayār āhe.}
\end{quote}

\begin{quote}
The first draft of the letter is ready.
\end{quote}

Masculine.

\begin{cousinweb}
\textbf{\mr{मसुदा}} (\emph{masudā}) is a loan through Persian usually traced to Arabic \emph{musawwada}, \textquotedblleft{}a rough copy\textquotedblright{}.
\end{cousinweb}

\begin{grammarlens}[title={What you've built}]
One more word for saying what happens if.
\end{grammarlens}

\subsection*{Your turn}
\begin{itemize}
\raggedright
  \item \emph{Say it:} \emph{masudā}
  \item \emph{Say it:} \emph{masudā}, once more
  \item \emph{Say it:} \emph{patrācā pahilā masudā tayār āhe.} — The first draft of the letter is ready.
  \item \emph{Recall:} say the Marathi for an identity card, then the Marathi for an outline, then say \emph{masudā} again
\end{itemize}

\begin{tcolorbox}[breakable,colback=teal!4,colframe=teal!35!black,title={Before you move on}]
What does \mr{मसुदा} mean? (\textquotedblleft{}A draft\textquotedblright{}.) Say it once more.
\end{tcolorbox}
\section[\texorpdfstring{\mr{नियोजन} \mr{करणे} (niyojan karṇe)}{नियोजन करणे (niyojan karṇe)}]{\mr{नियोजन} \mr{करणे} (niyojan karṇe) — to plan ahead}

\label{lesson:MR-C544-niyojan-karne}



\begin{quote}
Before the new one: say the Marathi for an outline, then the Marathi for a draft.
\end{quote}

\subsection*{You'll want to know: \mr{नियोजन} \mr{करणे}}
\textbf{\mr{नियोजन} \mr{करणे}} — \emph{niyojan karṇe} — \textquotedblleft{}to plan ahead\textquotedblright{}.

\begin{quote}
\mr{सहलीसाठी} \mr{आम्ही} \mr{आधीच} \mr{सगळ्याचं} \mr{नियोजन} \mr{केलं}.
\end{quote}

\begin{quote}
\emph{sahlīsāṭhī āmhī ādhīc sagḷyācaṁ niyojan kelaṁ.}
\end{quote}

\begin{quote}
We planned everything for the trip in advance.
\end{quote}

Past, with -\mr{ने} (he …): \mr{त्याने} \mr{नियोजन} \mr{केलं}, \emph{tyāne niyojan kelaṁ}.

What is planned takes -\mr{चं}: \mr{खर्चाचं} \mr{नियोजन} \mr{करणे}.

\begin{cousinweb}
\textbf{\mr{नियोजन} \mr{करणे}} (\emph{niyojan karṇe}) is the tatsama \mr{नियोजन}, from Sanskrit \emph{niyojana}, \textquotedblleft{}fastening, setting to a task\textquotedblright{} (Monier-Williams), with \mr{करणे}.
\end{cousinweb}

\begin{grammarlens}[title={What you've built}]
One more word for saying what happens if.
\end{grammarlens}

\subsection*{Your turn}
\begin{itemize}
\raggedright
  \item \emph{Say it:} \emph{niyojan karṇe}
  \item \emph{Say it:} \emph{niyojan karṇe}, once more
  \item \emph{Say it:} \emph{sahlīsāṭhī āmhī ādhīc sagḷyācaṁ niyojan kelaṁ.} — We planned everything for the trip in advance.
  \item \emph{Recall:} say the Marathi for an outline, then the Marathi for a draft, then say \emph{niyojan karṇe} again
\end{itemize}

\begin{tcolorbox}[breakable,colback=teal!4,colframe=teal!35!black,title={Before you move on}]
What does \mr{नियोजन} \mr{करणे} mean? (\textquotedblleft{}To plan ahead\textquotedblright{}.) Say it once more.
\end{tcolorbox}
\section[\texorpdfstring{\mr{लांबणीवर} \mr{टाकणे} (lāmbṇīvar ṭākṇe)}{लांबणीवर टाकणे (lāmbṇīvar ṭākṇe)}]{\mr{लांबणीवर} \mr{टाकणे} (lāmbṇīvar ṭākṇe) — to postpone}

\label{lesson:MR-C544-lambnivar-takne}



\begin{quote}
Before the new one: say the Marathi for a draft, then the Marathi for to plan ahead.
\end{quote}

\subsection*{You'll want to know: \mr{लांबणीवर} \mr{टाकणे}}
\textbf{\mr{लांबणीवर} \mr{टाकणे}} — \emph{lāmbṇīvar ṭākṇe} — \textquotedblleft{}to postpone\textquotedblright{}.

\begin{quote}
\mr{पावसामुळे} \mr{पंचांनी} \mr{सामना} \mr{लांबणीवर} \mr{टाकला}.
\end{quote}

\begin{quote}
\emph{pāvsāmuḷe pañcānnī sāmnā lāmbṇīvar ṭāklā.}
\end{quote}

\begin{quote}
Because of the rain the umpires postponed the match.
\end{quote}

Past, with -\mr{ने} (he …): \mr{त्याने} \mr{बैठक} \mr{लांबणीवर} \mr{टाकली}, \emph{tyāne baiṭhak lāmbṇīvar ṭāklī}.

Literally to put on the long side.

\begin{cousinweb}
\textbf{\mr{लांबणीवर} \mr{टाकणे}} (\emph{lāmbṇīvar ṭākṇe}) is built on \mr{लांबणी}, \textquotedblleft{}lengthening\textquotedblright{}, from \mr{लांब}, \textquotedblleft{}long\textquotedblright{}, inherited from Sanskrit \emph{lamba}, \textquotedblleft{}hanging down, long\textquotedblright{} (CDIAL 10951), with \mr{टाकणे}, \textquotedblleft{}to put\textquotedblright{}.
\end{cousinweb}

\begin{grammarlens}[title={What you've built}]
One more word for saying what happens if.
\end{grammarlens}

\subsection*{Your turn}
\begin{itemize}
\raggedright
  \item \emph{Say it:} \emph{lāmbṇīvar ṭākṇe}
  \item \emph{Say it:} \emph{lāmbṇīvar ṭākṇe}, once more
  \item \emph{Say it:} \emph{pāvsāmuḷe pañcānnī sāmnā lāmbṇīvar ṭāklā.} — Because of the rain the umpires postponed the match.
  \item \emph{Recall:} say the Marathi for a draft, then the Marathi for to plan ahead, then say \emph{lāmbṇīvar ṭākṇe} again
\end{itemize}

\begin{tcolorbox}[breakable,colback=teal!4,colframe=teal!35!black,title={Before you move on}]
What does \mr{लांबणीवर} \mr{टाकणे} mean? (\textquotedblleft{}To postpone\textquotedblright{}.) Say it once more.
\end{tcolorbox}
\section[\texorpdfstring{\mr{तडीस} \mr{नेणे} (taḍīs neṇe)}{तडीस नेणे (taḍīs neṇe)}]{\mr{तडीस} \mr{नेणे} (taḍīs neṇe) — to see through; to complete}

\label{lesson:MR-C544-tadis-nene}



\begin{quote}
Before the new one: say the Marathi for to plan ahead, then the Marathi for to postpone.
\end{quote}

\subsection*{You'll want to know: \mr{तडीस} \mr{नेणे}}
\textbf{\mr{तडीस} \mr{नेणे}} — \emph{taḍīs neṇe} — \textquotedblleft{}to see through; to complete\textquotedblright{}.

\begin{quote}
\mr{पैसे} \mr{कमी} \mr{होते}, \mr{तरी} \mr{त्याने} \mr{योजना} \mr{तडीस} \mr{नेली}.
\end{quote}

\begin{quote}
\emph{paise kamī hote, tarī tyāne yojnā taḍīs nelī.}
\end{quote}

\begin{quote}
Money was short, but he saw the plan through.
\end{quote}

Past, with -\mr{ने} (he …): \mr{त्याने} \mr{काम} \mr{तडीस} \mr{नेलं}, \emph{tyāne kām taḍīs nelaṁ}.

Literally to take it to the shore: \mr{तड} is a bank.

\begin{cousinweb}
\textbf{\mr{तडीस} \mr{नेणे}} (\emph{taḍīs neṇe}) is \mr{तडीस}, \textquotedblleft{}to the bank\textquotedblright{}, built on \mr{तड}, which is inherited from Sanskrit \emph{taṭa}, \textquotedblleft{}slope, shore\textquotedblright{} (CDIAL 5629), with \mr{नेणे}, \textquotedblleft{}to take\textquotedblright{}.
\end{cousinweb}

\begin{grammarlens}[title={What you've built}]
One more word for saying what happens if.
\end{grammarlens}

\subsection*{Your turn}
\begin{itemize}
\raggedright
  \item \emph{Say it:} \emph{taḍīs neṇe}
  \item \emph{Say it:} \emph{taḍīs neṇe}, once more
  \item \emph{Say it:} \emph{paise kamī hote, tarī tyāne yojnā taḍīs nelī.} — Money was short, but he saw the plan through.
  \item \emph{Recall:} say the Marathi for to plan ahead, then the Marathi for to postpone, then say \emph{taḍīs neṇe} again
\end{itemize}

\begin{tcolorbox}[breakable,colback=teal!4,colframe=teal!35!black,title={Before you move on}]
What does \mr{तडीस} \mr{नेणे} mean? (\textquotedblleft{}To see through; to complete\textquotedblright{}.) Say it once more.
\end{tcolorbox}
